\begin{textsample}{POS Dim 3 – mg – Score 12.00 – mg/mg\_inf\_organizacao\_2.txt}  \label{ex:f3_pos_247}
2.5.1 Direitos de Aprendizagem A BNCC ( 2017 ) preconiza seis direitos de aprendizagem que devem ser garantidos na Educação \textbf{Infantil}, considerando as diferentes maneiras pelas quais \textbf{bebês} e \textbf{crianças} aprendem e constroem sentidos sobre si, os outros e o mundo ; as especificidades e as características da vida contemporânea e a incorporação da Educação \textbf{Infantil} no sistema educacional.
Esses direitos são : - CONVIVER com outras \textbf{crianças} e \textbf{adultos}, em \textbf{pequenos} e grandes grupos, utilizando diferentes linguagens, ampliando o conhecimento de si e do outro, o respeito em relação à cultura e às diferenças entre as pessoas. - \textbf{BRINCAR} de diversas formas, em diferentes espaços e tempos, com diferentes parceiros ( \textbf{crianças} e \textbf{adultos} ), de forma a ampliar e diversificar suas possibilidades de acesso a produções culturais.
A participação e as transformações introduzidas pelas \textbf{crianças} nas \textbf{brincadeiras} devem ser valorizadas, tendo em vista o estímulo ao desenvolvimento de seus conhecimentos, sua imaginação, criatividade, experiências emocionais, corporais, \textbf{sensoriais}, expressivas, cognitivas, sociais e relacionais. - PARTICIPAR ativamente, com \textbf{adultos} e outras \textbf{crianças}, tanto do planejamento da gestão da escola e das atividades propostas pelo educador quanto da realização das atividades da vida cotidiana, tais como a escolha das \textbf{brincadeiras}, dos materiais e dos ambientes, desenvolvendo diferentes linguagens e elaborando conhecimentos, decidindo e se posicionando. - EXPLORAR movimentos, \textbf{gestos}, \textbf{sons}, formas, \textbf{texturas}, cores, palavras, emoções, transformações, relacionamentos, histórias, objetos, elementos da natureza, na escola e fora dela, ampliando seus saberes sobre a cultura, em suas diversas modalidades : as artes, a escrita, a ciência e a tecnologia. - EXPRESSAR, como sujeito dialógico, criativo e sensível, suas necessidades, emoções, sentimentos, dúvidas, hipóteses, descobertas, opiniões, questionamentos, por meio de diferentes linguagens. - CONHECER -SE e construir sua identidade pessoal, social e cultural, constituindo uma imagem positiva de si e de seus grupos de pertencimento, nas diversas experiências de \textbf{cuidados}, interações, \textbf{brincadeiras}, linguagens vivenciadas na instituição escolar, em seu contexto familiar e comunitário.
Para que estes direitos sejam contemplados, a BNCC estruturou as aprendizagens a serem alcançadas em toda etapa da Educação \textbf{Infantil} entrelaçando as situações e as experiências cotidianas da \textbf{criança} e seus saberes, aos conhecimentos que fazem parte do patrimônio cultural, artístico, científico e tecnológico, por meio dos campos de experiências.
Em cada campo de experiências os direitos de aprendizagem ganham especificidades, dessa forma, foram reelaborados, traduzindo aspectos peculiares, conforme estes se apresentam.

% matched lemmas: adulto, bebê, brincadeira, brincar, criança, cuidado, gesto, infantil, pequeno, sensorial, som, textura
\end{textsample}
